\subsection{Conversational Agents in Healthcare}

Conversational agents in healthcare have been studied as tools for service support, patient education, coaching, triage, monitoring, and behavioural interventions. Tudor Car \textit{et al.} (2020) reported that the healthcare chatbot literature remained diverse in purpose and format, but often lacked robust evidence for safety, effectiveness, and acceptability. They also observed that many systems were text-based, AI-driven, and mobile or web delivered, which aligned with the deployment assumptions of the present project.

Milne-Ives \textit{et al.} (2020) reviewed the effectiveness of AI conversational agents in healthcare and found encouraging use across multiple tasks, but also highlighted methodological variation and the need for stronger evaluation. Denecke and May (2022) focused specifically on usability assessment and concluded that no standardised procedure yet existed for evaluating healthcare conversational agents. This was important for the current project because it suggested that usability and operational traceability should be documented explicitly rather than assumed.

A recurring conclusion in the literature was that conversational agents should support rather than replace healthcare professionals (Tudor Car \textit{et al.}, 2020). That conclusion directly informed the scope of the chatbot developed here. The project therefore avoided any claim that the system could act as a clinician. Instead, it was positioned as a bounded laboratory assistant with controlled informational output.

\subsection{Sequential Models for Intent Classification}

Intent classification in text-based systems is a supervised sequence classification problem. The order of tokens matters because meaning depends on context, phrase structure, and local dependencies. Recurrent neural networks were developed to model sequential data, but classic RNNs suffered from vanishing and exploding gradients over longer sequences. Long Short-Term Memory was introduced to address that problem through gated memory cells capable of preserving relevant state over time (Hochreiter and Schmidhuber, 1997).

Bidirectional LSTM extended this idea by processing the sequence in forward and backward directions, thereby allowing the representation at a given position to incorporate both left and right context. Graves and Schmidhuber (2005) demonstrated that bidirectional LSTM was effective in sequence tasks where contextual information was crucial. Although the cited work addressed phoneme classification rather than text intent classification, the architectural principle remained relevant: bidirectional processing is useful where interpretation depends on surrounding terms.

The project deliberately chose BiLSTM rather than a transformer-based encoder. That decision was taken for practical and methodological reasons. The available dataset was modest rather than large scale; the deployment target required efficient local inference; and the final-year project title explicitly foregrounded sequential models. In addition, a sequential model with explicit vocabulary, token IDs, and fixed padding made the inference pipeline easier to inspect in the admin trace interface. This supported the educational and evaluative aims of the project.

The limitation of this choice was also acknowledged. Transformer models typically generalise more effectively across varied semantic phrasing when substantial pretraining is available. Consequently, the sequential model required more careful intent boundary design, broader paraphrase coverage, and a stronger data expansion process.

At the NLP pipeline level, sequential modelling was only one part of the design. The implemented system also used rule-based preprocessing, lexical tokenisation, vocabulary encoding, lightweight language detection, rule-based intent short-circuiting for safety or conversational cases, domain-specific entity detection, and later-stage information extraction for report values. This distinction was important for the dissertation because it showed that the artefact was not a pure end-to-end neural chatbot. Instead, it was a layered NLP system in which each technique solved a specific part of the problem.

\subsection{Retrieval-Based Response Generation}

Response generation in task-specific medical assistants can be approached in several ways: hand-written responses, rule-based templates, generative neural models, or retrieval-based methods. A purely rule-based chatbot often becomes brittle because language variation rapidly exceeds manually encoded patterns. A free-text generative model may be more flexible, but the medical domain introduces safety risks, reduced controllability, and greater difficulty in validating responses.

For that reason, a retrieval-based strategy was adopted. The assistant first determined what kind of question had been asked, then selected an answer from a curated knowledge base. This approach aligned with information retrieval principles that treat documents or question-answer entries as weighted term vectors. Salton and Buckley (1988) showed that term-weighting schemes based on weighted single terms were highly effective for automatic text retrieval. In the present system, TF-IDF retrieval was used because it was simple, interpretable, efficient on small corpora, and compatible with controlled answer sets.

The retrieval-based layer was not designed as a replacement for classification. Instead, it complemented the classifier. The intent model determined the response space, while retrieval selected the most appropriate answer within that space. This division of responsibility was important. When the classifier had already identified a narrow operational intent such as \texttt{sample\_collection} or \texttt{quality\_control}, broad entity overrides were not allowed to replace that intent with generic CBC explanations. This design rule reduced wrong-but-plausible responses.

\subsection{MLOps, Monitoring, and Technical Debt}

MLOps extends DevOps principles into the machine learning lifecycle. Kreuzberger, Kühl and Hirschl (2023) described MLOps as involving principles, workflows, roles, and technical components across data collection, validation, training, deployment, monitoring, and retraining. That view was strongly relevant to this project because the chatbot was not merely trained once and frozen. Instead, it was designed to ingest new labelled phrases, rebuild datasets, regenerate fixed splits, retrain specific model profiles, and redeploy with updated artefacts.

Sculley \textit{et al.} (2015) argued that machine learning systems incur hidden technical debt through entanglement, undeclared consumers, data dependencies, and configuration issues. The present project reflected that warning directly. Separate version-relevant artefacts were maintained for configuration, labelled data, derived datasets, fixed splits, model weights, evaluation runs, and knowledge-base content. The admin panel was then used to expose those artefacts through a practical monitoring interface. That design did not eliminate technical debt entirely, but it reduced the invisibility of system state.

\subsection{Research Gap and Justification}

The literature showed a clear need for healthcare chatbots that were safe, scoped, and evaluable. It also showed that MLOps maturity was important for real-world ML systems. However, there remained a gap between those discussions and small-scale, domain-specific educational projects. Many student chatbot systems demonstrate inference, but do not include data versioning, logging, retraining, evaluation discipline, or operational review tools. Conversely, many MLOps examples assume large cloud platforms and do not focus on a bounded medical assistant.

The present project therefore occupied a useful middle ground. It did not claim clinical-scale innovation. Instead, it implemented a realistic and inspectable architecture that connected sequential NLP, retrieval-based generation, deployment controls, and review-driven retraining within a manageable final-year project scope.
